\section{Chapters 18 and 19}

\paragraph{Chapter 18}

\begin{verse}
The Way was lost\\ ``Benevolence and righteousness'' were put in its place\\ By skill and cunning (having lost original simplicity)\\ Great hypocrisy ensued\\ The chain of blood that was tied to the origin was broken\\ And was replaced with family sentiments\\ When the kingdoms fell into disorder\\ ``Good ministers'' became prominent.
\end{verse}

\paragraph{Chapter 19}

\begin{verse}
Remove (petty) wisdom, put prudence aside\\ And the advantage to the people will be very great\\ Humanitarianism and morality are not proclaimed\\ And a spontaneous solidarity = from the family will return to rule in society\\ Utilitarian goals will be devalued, self-absorbed action will be despised\\ And guilty acts will vanish in society.\\ (But even) these three particularly\\ Do not have to be followed to the letter\\ What is essentially necessary:\\ To be (oneself) in rough, natural, candour\\ Detached from particular ambitions, liberated from vain desires.
\end{verse}

\paragraph{Commentary}


There is a visible continuity between these two chapters. Proclaiming ``virtues'' is for Taoism merely the sign of a corrupt society. See Chapter 38: ``Superior virtues do not announce themselves as virtues.'' It speaks instead of the ``non-virtues of virtues'', the latter understood as artificial corresponding to a pedantic exterior that, more than restoring an integral life, will only increase the damage by favouring hypocrisy and cunning. The Taoist argument in that sense, not without reference to a second-rate or poorly understood Confucianism, often went beyond the sign. ``Benevolence and righteousness'' --- \emph{ren} and \emph{yi} --- are precisely the cardinal Confucian virtues. Even family sentiments reduced to conventional and obligatory behaviour are decadent. (Chapter 18) See \textbf{Chuang Tzu} (XIII,5):

\begin{quotex}
These concerns about benevolence and righteousness recall a man who beats a drum while searching for a fleeing son, with the only effect of making him flee even further... Unite your influence to that of the Principle rather than imposing artificial virtues, and you will be able to achieve something.
\end{quotex}


From the last line of Chapter 19 the positive reference point is clarified: conduct derives from being oneself, rather than reflecting a norm of society. The elimination of desires, as vain desires, must not be understood here in an ascetic, forced, or renunciatory sense. It is the fundamental idea of Taoism, that in man a complex of desires, tendencies and interests is not at all ``natural''; it is an error, it is something artificial and parasitic. One is truly oneself --- and one is on the Path --- when one is liberated from it. Then true virtues are also manifested. ``Morality'', far from restoring this state of natural equilibrium, moves even further from it, adding another stratification to a stratification . An effective image:

\begin{quotex}
The beginning is what I call the thaw (the dissolution of the concretion of the exterior I), after that the stream begins to take its course. (Chuang Tzu, XXIII, 3)
\end{quotex}


In a certain way, the same double application is reaffirmed in the chapter that follows. Among the excrescences in respect to the natural state, which one must prune, the world of intellectualism and of the social distinctions between ``good'' and ``evil'' is also included. However, in contrast, being oneself is what brings you quite far, along the line of the Transcendent Man in his closed and impenetrable form. The reference, as will be seen, is in the first person --- it is almost the only passage of the \emph{Tao Te Ching} in which the author speaks of himself: however, with the visible intent of sketching out a model.

\paragraph{Chinese text and literal translation}


Chapter 18 (\begin{CJK*}{UTF8}{bsmi}第十八章\end{CJK*})
\columnratio{.37}
\begin{paracol}{2}
\begin{CJK*}{UTF8}{bsmi}
\begin{verse}
大道廢，有仁義；\\ 智慧出，有大偽；\\ 六親不和有孝慈，\\ 國家昏亂有忠臣。
\end{verse}
\end{CJK*}
\switchcolumn
\begin{verse}
When the Dao is lost, so there arises benevolence and righteousness.\\ When prudence and wisdom appear, there is great hypocrisy.\\ When family relations are not right, there is filial piety and paternal affection.\\ When the state is confused and chaotic, there is loyalty and faithfulness.
\end{verse}
\end{paracol}


Chapter 19 (\begin{CJK*}{UTF8}{bsmi}第十九章\end{CJK*})
\columnratio{.39}
\begin{paracol}{2}
\begin{verse}
\begin{CJK*}{UTF8}{bsmi}
絕聖棄智，民利百倍；\\ 絕仁棄義，民復孝慈；\\ 絕巧棄利，盜賊無有；\\ 此三者，以為文不足，\\ 故令有所屬：\\ 見素抱樸，\\ 少私寡欲。
\end{CJK*}
\end{verse}
\switchcolumn
\vspace*{-.2em}
\begin{verse}
Abandon holiness, relinquish prudence; the people will flourish a hundredfold,\\ Abandon benevolence, relinquish righteousness; and people will return to filial piety and affection.\\ Abandon cleverness, relinquish gain; and thieves and robbers will not appear.\\ As I know such three to be no mere words,\\ Stay true to that which is reliable.\\ Recognise simplicity, embrace purity, \\ Lessen the self, diminish desires.
\end{verse}
\end{paracol}

\flrightit{Posted on 2021-03-04 by Lao Tzu}

